%! Least-Squares Regression — Unit 2, Lesson 2.4 — Homework (Answer Key)
% GRADUAL RELEASE: this is the lesson's SCORED practice and the LAST component in the packet.
% Paper homework is the DEFAULT for this course; DeltaMath is an occasional override that
% replaces this page and strikes its score cell on the cover. Started in the last eight minutes
% of the period and finished at home. Due the first class after two study halls.
% THIRD CONTEXT: the notes teach on sleep and reaction time, Guided Practice runs on used-car
% prices, and these items are on car weight and highway mileage. ONE PAGE, five items.
% THE DATA: 35 car models, weight x (1000s of lb, 2.5 to 5.0), highway mpg y.
%   xbar 3.5, s_x 0.6, ybar 30, s_y 5, r -0.84 -> b = -7.0, a = 54.5, r^2 = 0.7056.
% \answerspace{H}{} reserves exactly H in the blank; the key passes the answer as the second
% argument, so the two files paginate identically by construction. Keep the macro and every
% \boxguard byte-identical in the blank and the key.
% GENERATED FROM homework/main.tex — the blank with answers filled in.
% NO teachernote here — scoring guidance lives in the lesson plan.
\documentclass[12pt]{article}
\usepackage{apstats-article}
\usepackage{apstats-key}   % pulls in -boxes; do NOT also load -boxes
\usepackage{grading}
\gradingassignment{apstats-u02-l04-hw}
\gradingcourse{advanced_statistics}
\newcommand{\answerspace}[2]{\par\nopagebreak\noindent\begin{minipage}[t][#1][t]{\linewidth}%
  \color{keyred}\bfseries #2\end{minipage}\par}

\begin{document}

\pageheader{Unit 2, Lesson 2.4}{Homework --- Answer Key}
% NAMESTRIP: no \namedateperiod here — the name row lives on the cover only.

\vspace{0.04in}

% Authored BYTE-IDENTICALLY in homework/ and homework_key/.
\boxguard[8]
\begin{remindbox}
\small
\textbf{This is your graded homework.} It is scored and due the first class after two study
halls. Start it now; ask about anything that stalls you before you leave.
\end{remindbox}

\vspace{0.04in}

\boxguard[14]
\begin{notesbox}{Weight and Highway Mileage}
\small
For 35 car models: \textbf{weight} ($x$, 1000s of pounds, 2.5 to 5.0) and \textbf{highway
mileage} ($y$, mpg). \ $\bar x = 3.5$, \ $s_x = 0.6$, \ $\bar y = 30$, \ $s_y = 5$, \ $r = -0.84$. \hfill \textbf{Part A --- Build the line.}
\begin{enumerate}[label=\textbf{A\arabic*.}, leftmargin=2.2em, itemsep=2pt]
  \item Find $b$ and $a$ (show your work), then write the LSRL with the variable names.
        \answerspace{1.5cm}{$b = -0.84 \cdot 5/0.6 = -7.0$ mpg per 1,000 lb; \ $a = 30 - (-7.0)(3.5) = 54.5$ mpg.\par $\widehat{\text{mpg}} = 54.5 - 7.0(\text{weight})$\par\vfill\hfill\normalcolor\mdseries\textit{b:}~\gradedblank[id=hw1a, type=short, points=1, answer=-7.0, check=numeric, tolerance=0.05]{0.65in}{0.24in}\quad\textit{a:}~\gradedblank[id=hw1b, type=short, points=1, answer=54.5, check=numeric, tolerance=0.1]{0.65in}{0.24in}\quad\textit{LSRL:}~\gradedblank[id=hw1c, type=rubric, points=1, display={$\widehat{\text{mpg}} = 54.5 - 7.0(\text{weight})$}]{1.9in}{0.24in}}
  \item Interpret the slope in context.
        \answerspace{1.15cm}{\normalcolor\mdseries\hspace*{10.5pt}\gradedblank[id=hw2, type=rubric, points=2, display={For each additional 1,000 pounds of weight, the \emph{predicted} highway mileage drops by 7.0 mpg.}]{\linewidth-10.5pt}{0.95cm}}
  \item Interpret the intercept in context. Does it have a sensible meaning here? Explain.
        \answerspace{1.15cm}{\normalcolor\mdseries\hspace*{10.5pt}\gradedblank[id=hw3, type=rubric, points=2, display={Predicted 54.5 mpg for a car weighing 0 pounds. Not meaningful: no car weighs 0 lb, far outside 2.5--5.0.}]{\linewidth-10.5pt}{0.95cm}}
\end{enumerate}

\textbf{Part B --- How good is it?}
\begin{enumerate}[label=\textbf{B\arabic*.}, leftmargin=2.2em, itemsep=2pt]
  \item Find $r^2$ and interpret it in context.
        \answerspace{1.15cm}{\normalcolor\mdseries\hspace*{10.5pt}\gradedblank[id=hw4b, type=rubric, points=2, display={About 71\% of the variation in highway mileage is explained by the linear model with weight.}]{\linewidth-10.5pt-1.0in-6pt}{0.95cm}\hfill\gradedblank[id=hw4a, type=short, points=1, answer=0.71, check=numeric, tolerance=0.01]{1.0in}{0.85cm}}
  \item \textbf{AP-style.} Which statement is \textbf{correct}? Circle one, then justify it in a
        sentence.
        \begin{enumerate}[label=\textbf{(\Alph*)}, leftmargin=2.2em, itemsep=0pt]
          \item Every extra 1,000 pounds lowers a car's mileage by exactly 7 mpg.
          \item \ans{For each extra 1,000 pounds, the predicted highway mileage drops by 7 mpg.}
          \item About 84\% of the variation in mileage is explained by weight.
          \item A car weighing 0 pounds would get about 54.5 mpg.
          \item Adding weight to a car causes its mileage to fall by 7 mpg.
                \hfill\gradedright{\textit{Letter:}~\gradedblank[id=hw5, type=choice, points=1, answer=B]{0.45in}{0.24in}}
        \end{enumerate}
        \answerspace{0.8cm}{\normalcolor\mdseries\hspace*{10.5pt}\gradedblank[id=hw5j, type=rubric, points=1, display={(B): the slope gives the change in \emph{predicted} mileage.}]{\linewidth-10.5pt}{0.7cm}}
\end{enumerate}
\end{notesbox}

\vspace{0.04in}

% The next-lesson preview closes the packet, so it lives here rather than in AP Practice.
% A one-line, unbreakable box: guard it by its own height so it moves whole.
\boxguard[3]
\begin{spiralbox}
\small
\textbf{Next: Departures from Linearity} --- what one unusual point can do to $b$, $a$, and $r^2$.
\end{spiralbox}

\end{document}
